\documentclass[11pt, reqno]{article}

\usepackage{amsmath, amsthm, amssymb}
\usepackage{enumitem}
\usepackage{tcolorbox}
\usepackage{hyperref}
\usepackage{tikz}
\usepackage{tikz-cd}
\usepackage{pgfplots}
\pgfplotsset{compat=1.18}
\usetikzlibrary{arrows.meta}
\usepackage{mathrsfs}
\usepackage{fancyhdr}
\usepackage[bottom=0.75in, top=1in, left=0.5in, right=0.5in]{geometry}
\usepackage{array}   % for \newcolumntype macro
\newcolumntype{L}{>{$}l<{$}}

\theoremstyle{plain}
\newtheorem*{theorem}{Theorem}
\newtheorem*{proposition}{Proposition}
\newtheorem{exercise}{Exercise}
\newtheorem*{lemma}{Lemma}
\newtheorem*{corollary}{Corollary}

\theoremstyle{definition}
\newtheorem*{definition}{Definition}
\newtheorem*{example}{Example}

\theoremstyle{remark}
\newtheorem*{remark}{Remark}

\renewcommand{\phi}{\varphi}
\renewcommand{\epsilon}{\varepsilon}
\renewcommand{\emptyset}{\varnothing}

\newcommand{\RR}{\mathbb{R}}
\newcommand{\ZZ}{\mathbb{Z}}
\newcommand{\NN}{\mathbb{N}}
\newcommand{\CC}{\mathbb{C}}
\newcommand{\QQ}{\mathbb{Q}}

\DeclareMathOperator{\ima}{\text{im}}

\begin{document}

\topmargin=-40pt
\rhead{Henry Woodburn}
\lhead{Math 658}
\renewcommand{\headrulewidth}{1pt}
\renewcommand{\headsep}{20pt}
\thispagestyle{fancy}

{\Huge \bfseries \noindent Homework 2}

\begin{enumerate}
    \item[1.] (a.) The derivative of a distribution $f \in \mathscr{S}'$ is defined by the rule
    \[
        \langle \partial f, \phi\rangle = - \langle f, \partial \phi\rangle,\quad \text{for}\ \phi \in \mathscr{S}.
    \]
    To compute the derivative of $\chi_{[a,b]}$, letting $\phi \in \mathscr{S}$, we have
    \[
        \langle \partial_x \chi_{[a,b]}, \phi\rangle = -\langle \chi_{[a,b]}, \partial_x \phi\rangle = -\int_a^b \partial_x \phi = \phi(a) - \phi(b).
    \]
    So the action of $\partial_x \chi_{[a,b]}$ is the same as $\delta_a - \delta_b$, so this is the distributional derivative.

    (b.) Take $x_j = x$ without loss of generality. Again letting $\phi \in \mathscr{S}(\RR^2)$, we can write
    \begin{align*}
        \langle \partial_x \chi_{B(0,1)}, \phi\rangle &= -\int_{B(0,1)} \partial_x \phi(x, y) dV = 
        -\int_{-1}^1 \int_{-\sqrt{1 - y^2}}^{\sqrt{1 - y^2}}\partial_x \phi dx dy\\
        &= \int_{\partial B(0,1)^-} \phi dS - \int_{\partial B(0,1)^+} \phi dS,
    \end{align*}

    where $dS$ denotes a measure on the boundary $\partial B(0,1)$ and $-$ and $+$ denote left and right semicircles. This is the difference between the integral of $\phi$ over the left and
    right semicircles.

    (c.) We compute the fourier transforms of $\sin(x)$ and $\cos(x)$, which are not integrable functions in the $L^1$ sense. We use the definition
    \[
        \langle \mathscr{F}u, \phi\rangle = \langle u, \mathscr{F}\phi\rangle,\quad \text{for}\ \phi \in \mathscr{S}.
    \]
    We already know that $\hat{\left(e^{2i \pi x \cdot a}\right)} = \delta_a$ for $a \in \RR$, so that $\hat{\left(e^{ix}\right)} = \delta_{\frac{1}{2\pi}}$. 
    Moreover, $\sin(x)$ can be expressed as
    \[
        \sin(x) = \frac{e^{ix} - e^{-ix}}{2i}
    \]
    so that by linearity of the fourier transform, we have
    \[
        \mathscr{F}(\sin(x)) = \frac{1}{2i}\left(\delta_{\frac{1}{2\pi}} - \delta_{\frac{-1}{2\pi}}\right).
    \]

    (d.) We need to compute
    \[
        \langle \partial \log|x|, \phi\rangle = -\langle \log|x|,\partial_x \phi\rangle
    \]
    for any $\phi \in \mathscr{S}(\RR)$. Let $\epsilon > 0$. Then
    \[
        \int_{\epsilon \leq |x|} \phi(x) \frac{1}{x}dx = -\int_{\epsilon \leq |x|}\partial_x\phi(x) \log|x|dx
    \]
    after integrating by parts, noting that the first integrand is integrable for any $\epsilon > 0$. 
    Moreover, the function $\log|x|$ is integrable on $\RR$, so we can let epsilon go to zero in the final integral.
    So the limit converges for all $\phi$:
    \[
        u(\phi) = \lim_{\epsilon \to 0} \int_{\epsilon \leq |x|} \phi(x)\frac{1}{x}dx = -\int_\RR \partial \phi(x) \log|x|dx = \langle \log|x|, \partial_x \phi\rangle =: \langle \partial_x \log|x|, \phi\rangle.
    \]

    We have proven the sequence of tempered distributions defining $u$ converges to $\partial_x \log |x|$ in the weak$^*$ topology on $\mathscr{S}'(\RR)$, 
    and so this is the derivative.
    
\end{enumerate}

\end{document}